%-----------------------------------------------------------------------
\section{Threat Model and Security Assumptions}
\label{sec:threat_model}
%-----------------------------------------------------------------------

To contextualize the urgency and scope of IoT migration towards post-quantum standards, we formally define our assumed threat geometry governing resource-constrained interactions.

\subsection{Quantum Adversarial Model}

We adopt the standard Dolev-Yao model: the adversary $\mathcal{A}_{\text{CRQC}}$ can intercept, modify, and reroute any message on the IoT transmission medium (LoRaWAN, NB-IoT, Wi-Fi, or DSRC). This is unchanged from classical threat modeling.

What changes is computational power. $\mathcal{A}_{\text{CRQC}}$ has access to a cryptographically relevant quantum computer — one capable of running Shor's algorithm \cite{shor1994algorithms} against RSA/ECC key exchanges and Grover's algorithm \cite{grover1996fast} to halve the effective security of symmetric ciphers. We assume enough stable logical qubits to break any NIST Security Category 1 system ($\geq$ 128-bit classical equivalent) outright. Whether such a machine exists today is beside the point; the threat model must account for it because IoT firmware deployed now will still be running in 2040.

\subsection{The "Harvest Now, Decrypt Later" Immediacy}

A stolen RSA-encrypted packet is not just a breach today — it is a breach whenever a quantum computer arrives. Nation-state actors are already archiving bulk TLS traffic from IoT networks: medical telemetry, smart grid commands, industrial SCADA streams. If those records need to stay confidential for 25 years, and a fault-tolerant quantum computer appears within that window, every archived session is retroactively compromised. This "harvest now, decrypt later" timeline is what makes PQC migration urgent even before commercial quantum hardware matures.

\subsection{Side-Channel Vulnerabilities}

Beyond purely algebraic quantum breaches, $\mathcal{A}_{\text{CRQC}}$ is assumed to possess extensive localized hardware penetration metrics. Specifically:
\begin{itemize}
 \item \textbf{Timing Attacks:} Adversaries monitor minor CPU execution variations occurring during non-constant-time polynomial branch logic (e.g., Gaussian rejection sampling bounds typical in \falcon{}).
 \item \textbf{Differential Power Analysis (DPA):} Adversaries infer the secret lattice error term polynomial gradients strictly by correlating power trace consumption spikes executed by the IoT microcontroller mapping instructions to memory \cite{kocher1999differential, hutter2011side}.
\end{itemize}
These models heavily mandate that deployed IoT solutions execute strictly constant-time functions. Constant-time algorithms unconditionally execute at identical clock-cycle bounds regardless of input payloads, deliberately neutralizing external timing leakages at the explicit cost of substantial performance saturation.
